% Chapter 5 -- Troubleshooting (Mac OS 8 Installation Ch. 6 form; AN-14:
% bold symptom, then Cause / Solution).  The real messages, quoted from the
% source: os/boot/stage2.asm (the Non-System disk message, then a halt);
% os/milton/kmain.c (FLAIR-FAT-MOUNT-FAIL: the desktop starts without the
% data disk -- the disk icon has no name, the clock reads 12:00 am because
% the clock is read through InitechDOS; the PC LOAD LETTER lines "FLAIR
% heap exhausted" and "FLAIR needs a 640x480 framebuffer"; MSG-DOS-0014
% after a failed Exit to InitechDOS); os/milton/panic.c (PC LOAD LETTER
% (exception 0xhh), then the system halts); spec/dos_messages.json
% MSG-DOS-0014 / 0015; os/flair/event.c (one mouse button).
% GATES: test-flair-desktop test-flair-heap test-boot-system test-flair-exit-to-dos test-dosmsg test-pubs-flair-screens
\Gates{test-flair-desktop, test-flair-heap, test-boot-system,
  test-flair-exit-to-dos, test-dosmsg, test-pubs-flair-screens}
\chapter{Troubleshooting}
\index{troubleshooting}

\Symptom{Nothing appears on the screen}
\Cause The display is not turned on, or not connected.
\Solution Check that the display is plugged in, connected to the computer
and turned on, and that its brightness is turned up.
\Cause The display adapter cannot show 640 by 480 dots in 256 colours.
\Solution FLAIR needs a VGA adapter with the VESA BIOS Extensions,
Version~2.0 (Chapter~2).

\Symptom{The screen shows ``Non-System disk or disk error''}
The full message is:
\begin{TypeIn}
Non-System disk or disk error
Replace and strike any key when ready
\end{TypeIn}
\Cause The disk in the first drive is not the startup disk, or the startup
disk is damaged.
\Solution Put the startup disk in the first drive and restart the computer
with its reset button or by turning it off and on.  Pressing a key does not
restart it.  If the message appears again, use your other copy of the
startup disk.

\Symptom{The disk icon has no name, and the clock reads 12:00 am}
\Cause FLAIR started without the data disk.  The \texttt{INITECH} icon
stands for the second drive, but there is no disk to show; and FLAIR reads
the time through InitechDOS, which has no disk to work with.

\FlairScreen[scale=4, label=fig:gs-nodisk]{g_nodisk}{FLAIR started without
its data disk}
\Operation{FLAIR-FAT-MOUNT-FAIL}

\Solution Shut down, put the data disk in the second drive, and start the
computer again.

\Symptom{The screen shows ``PC LOAD LETTER''}
InitechOS has stopped because of a problem it cannot work around.  The
words in parentheses tell you which:
\begin{TypeIn}
PC LOAD LETTER  (FLAIR heap exhausted)
\end{TypeIn}
\Cause The computer has less memory than FLAIR needs.
\Solution FLAIR needs 4~MB of extended memory.  Have more memory
installed.
\begin{TypeIn}
PC LOAD LETTER  (FLAIR needs a 640x480 framebuffer)
\end{TypeIn}
\Cause The display adapter cannot give FLAIR a screen of 640 by 480 dots.
\Solution Use a display adapter with the VESA BIOS Extensions,
Version~2.0.
\par\noindent\hspace*{1.5em}\texttt{PC LOAD LETTER\ \ (exception 0x\textit{nn})}\par
\Cause A program has failed in a way the computer cannot recover from.  The
two characters after \texttt{0x} name the kind of failure.
\Solution Write down the whole message and restart the computer.  If the
message appears again when you repeat what you were doing, call the
Initech Help Desk (extension 2504) with the message you wrote down.

\Symptom{The pointer does not move}
\Cause The mouse is not connected.
\Solution Shut down if you can, turn off the computer, connect the mouse,
and turn the computer on again.

\Symptom{The pointer moves, but clicks do nothing}
\Cause You are pressing the right button of a two-button mouse.
\Solution FLAIR uses the left button only.

\Symptom{The screen shows ``Bad or missing Command Interpreter''}
\Cause You chose \menuitem{Exit to InitechDOS}, and InitechDOS could not
load its command processor from the startup disk.
\Solution Restart the computer.  If it happens again, make a new copy of
the startup disk from the original.

\Symptom{Other problems}
The \doctitle{FLAIR User's Guide}, Chapter~13, describes the problems you
may meet once FLAIR is running.
